\documentclass[11pt,a4paper]{article}

%% PACHETE MATEMATICE
\usepackage{amsmath,amssymb,amsfonts}
\usepackage{amsthm}
\usepackage{mathpazo}
\usepackage{eulervm}
\usepackage{enumerate}
\usepackage{color}
\usepackage{graphicx}
\usepackage[all]{xy}
\usepackage{xcolor}
\usepackage{framed}
\usepackage[unicode]{hyperref}
\setcounter{MaxMatrixCols}{20}
\usepackage{multicol}
%\usepackage[romanian]{babel}


%% ASPECT
\linespread{1.05}
\usepackage[T1]{fontenc}
\usepackage[utf8x]{inputenc}
\usepackage[margin=2cm]{geometry}
% \usepackage[color]{showkeys}
\usepackage{anyfontsize}
\usepackage{titlesec}
\titleformat*{\section}{\large\bfseries}

\usepackage{fancyhdr}
\pagestyle{fancy}
\rhead{\small \color{gray} Notițe de Adrian Manea}
\lhead{\small \color{gray} ETTI-AM2, 2017---2018, M. Joița \& A. Niță}


\usepackage[autostyle = false, style = german]{csquotes}
\usepackage[romanian]{babel}
\newcommand\qq{\enquote}

%% COMENZILE MELE
\renewcommand*{\proofname}{Dem.:}
\newcommand\bb{\mathbb}
\newcommand\dr{\mathrm}
\newcommand\kal{\mathcal}
\newcommand\fk{\mathfrak}
\newcommand\op{\oplus}
\newcommand\ot{\otimes}
\newcommand\ds{\displaystyle}
\newcommand\id{\indent}
\newcommand\nid{\noindent}
\newcommand\seq{\subseteq}
\newcommand\sneq{\subsetneq}
\newcommand\speq{\supseteq}
\newcommand\spneq{\supsetneq}
\newcommand\rar{\rightarrow}
\newcommand\Rar{\Rightarrow}
\newcommand\xrar{\xrightarrow}
\newcommand\lar{\leftarrow}
\newcommand\Lar{\Leftarrow}
\newcommand\xlar{\xleftarrow}
\newcommand\lrar{\leftrightarrow}
\newcommand\Lrar{\Leftrightarrow}
\newcommand\ideal{\trianglelefteq}
\newcommand\ai{\text{ a.\^{i}. }}
\newcommand\sk{(\textit{Schi\c{t}\u{a}}) }
\newcommand\rhup{\rightharpoonup}
\newcommand\rhdn{\rightharpoondown}
\newcommand\lhup{\leftharpoonup}
\newcommand\lhdn{\leftharpoondown}
\newcommand\vid{\emptyset}
\newcommand\wh{\widehat}
\newcommand\ol{\overline}
\newcommand\NN{\mathbb{N}}
\newcommand\RR{\mathbb{R}}
\newcommand\ZZ{\mathbb{Z}}
\newcommand\QQ{\mathbb{Q}}
\newcommand\CC{\mathbb{C}}

%% STILURILE MELE
\newtheoremstyle{dotless}{}{}{\itshape}{}{\bfseries}{:}{ }{}

\theoremstyle{dotless}
\newtheorem{theorem}{Teorem\u{a}}[section]
\newtheorem{proposition}{Propozi\c{t}ie}[section]
\newtheorem{lemma}{Lem\u{a}}[section]
\newtheorem{corollary}{Corolar}[section]
\newtheorem{exercise}{Exerci\c{t}iu}

\newtheoremstyle{dotlessdr}{}{}{}{}{\bfseries}{:}{ }{}
\theoremstyle{dotlessdr}
\newtheorem{example}{Exemplu}[section]
\newtheorem{definition}{Defini\c{t}ie}[section]
\newtheorem{remark}{Observa\c{t}ie}[section]




\begin{document}
% \definecolor{labelkey}{RGB}{222,0,0}

\begin{center}
  {\large\textbf{Seminar 2 \\
      Ecuații și sisteme diferențiale \\
      Exerciții suplimentare}}
\end{center}

\vspace{2cm}

1. Determinați soluția generală pentru următoarele ecuații diferențiale:
\begin{multicols}{2}
  \begin{enumerate}[(a)]
  \item $y' = 3y^{\frac{2}{3}}$;
  \item $xy' = 2y$;
  \item $y' = e^{x + y}$;
  \item $y' = \frac{2x}{3y^2 + 1}$;
  \item $(x^2 + y^2 + 2x) dx + 2xy dy = 0$;
  \item $(x^3 - 3xy^2 + 2)dx - (3x^2y - y^2) dy = 0$;
  \item $y^2 dx + (3^x - 2y) dy = 0$;
  \item $(x + y^2) dx - 2xy dy = 0$;
  \item $xy' + y = x^3$;
  \item $x^2 y' + (1 - 2x)y = x^2$;
  \end{enumerate}
\end{multicols}

\vspace{1cm}

2. Găsiți soluția generală pentru următoarele sisteme diferențiale:
\begin{enumerate}[(a)]
\item $\begin{cases}
    y' &= 2y + z \\
    z' &= y + 2z
  \end{cases}, \quad y = y(x), z = z(x)$;
\item $\begin{cases}
    y' &= -3y - 4z \\
    z' &= y + 2z
  \end{cases}, \quad y = y(x), z = z(x)$;
\item $\begin{cases}
    y_1' &= -2y_1 - y_2 + y_3 \\
    y_2' &= 5y_1 - y_2 + 4y_3 \\
    y_3' &= 5y_1 + y_2 + 2y_3
  \end{cases}, \quad y_1 = y_1(x), y_2 = y_2(x), y_3 = y_3(x)$;
\item $\begin{cases}
    x' &= x + 3y \\
    y' &= -x + 5y - 2e^t
  \end{cases}$, pentru funcțiile $x = x(t), y = y(t)$, cu condițiile $x(0) = 3, y(0) = 1$.
\end{enumerate}

\vspace{1cm}

3. Determinați soluția generală pentru următoarele ecuații diferențiale:
\begin{enumerate}[(a)]
\item $y^{(3)} + 4y^{(2)} + 3y' = 0$;
\item $y^{(2)} + 4y' + 4y = 0$;
\item $y^{(3)} + y = 0$;
\item $y^{(4)} + 4y^{(2)} = 0$.
\end{enumerate}

\vspace{1cm}

4. Determinați soluția generală pentru următoarele ecuații neomogene:
\begin{enumerate}[(a)]
\item $y^{(2)} - 2y' + y = \frac{1}{x} e^x$;
\item $y^{(2)} + y = \frac{1}{\cos x}$;
\item $y^{(3)} + y' = \tan x$;
\item $y^{(3)} + 2y^{(2)} = x + 2$;
\end{enumerate}

\vspace{1cm}

5. Determinați soluția sistemului diferențial liniar neomogen:
\[
  \begin{cases}
    x' &= 2x - y \\
    y' &= -x + 2y + e^t z \\
    z' &= -y + 2z
  \end{cases}
\]
știind că satisface condițiile inițiale $X(0) = \begin{pmatrix} 0 \\ 1 \\ 0 \end{pmatrix}$.

\vspace{1cm}

\textbf{Ecuații Euler}

O ecuație diferențială liniară de ordinul $n$ de forma:
\[
  L[y] = a_n x^n y^{(n)} + a_{n-1} x^{n-1} y^{(n-1)} + \dots + a_0 y = f(x)
\]
cu $a_i$ constante și $f \in \kal{C}^1(I)$ se numește \emph{ecuație Euler}.


Ecuațiile Euler se pot transforma în ecuații cu coeficienți constanți prin schimbarea
de variabilă $|x| = e^t$. De remarcat că, deoarece funcția necunoscută este $y = y(x)$,
pentru a obține $y'$ trebuie să aplicăm regula de derivare a funcțiilor compuse.
Fie $z(t) = y(e^t)$. Atunci avem, de exemplu, primele două derivate:
\begin{align*}
  z'(t) &= y'(e^t) e^t \Rightarrow y'(e^t) = e^{-t} z'(t) \\
  z^{\prime\prime}(t) &= y^{\prime\prime}(e^t) e^{2t} + y'(e^t)e^t =
                        y^{\prime\prime}(e^t) e^{2t} + z'(t)
\end{align*}


Exemplu: $x^2 y^{(2)} + xy' + y = x$.

Facem substituția $|x| = e^t$ și obținem:
\[
  e^{2t} \cdot e^{-2t}(y^{(2)} - y'(t)) - e^t \cdot y'(t) \cdot e^{-t} + y(t) = e^t.
\]
Echivalent:
\[
  y^{(2)}(t) - 2y'(t) + y(t) = e^t.
\]
Aceasta este o ecuație liniară de ordinul al doilea, neomogenă. Asociem ecuația
algebrică $f(r) = r^2 - 2r + 1 = (r - 1)^2 = 0$, deci soluția generală a ecuației
omogene este:
\[
  y(t) = c_1 e^t + c_2 te^t.
\]
Folosind apoi metoda variației constantelor, obținem succesiv:
\begin{align*}
  y(t) &= c_1 e^t + c_2te^t + \frac{t^2 e^t}{2} \Rightarrow \\
  \Rightarrow y(x) &= c_1 x + c_2 x \ln x + \frac{x\ln^2 x}{2}.
\end{align*}

\vspace{1cm}

6. Să se rezolve ecuațiile Euler:
\begin{enumerate}[(a)]
\item $x^2 y^{(2)} - 3xy' + 4y = 5x, \quad x > 0$;
\item $x^2 y^{(2)} - xy' + y = x + \ln x, \quad x > 0$;
\item $4(x + 1)^2 y^{(2)} + y = 0, \quad x > -1$.
\end{enumerate}



\end{document}

